\documentclass[10pt,a4paper]{amsart}
\usepackage[margin=19mm]{geometry}
\usepackage{amsmath,amssymb,mathtools}
\usepackage[scr=rsfs,cal=euler]{mathalfa}
\usepackage{xfrac}
\usepackage[unicode,hidelinks,bookmarks=false]{hyperref}
\newcommand{\ie}{i.e.\ }
\newcommand{\cf}{cf.\ }
\newcommand{\resp}{respectively\ }
\newcommand{\Subsection}{\subsection}
\providecommand{\nobreakdash}{\nobreak\hbox{-}}
\setlength{\emergencystretch}{2em}
\multlinegap0pt
\usepackage{tikz}
\usetikzlibrary{cd}
\usepackage{amssymb}
\usepackage{mathtools}
\usepackage{enumerate}
\usepackage{xcolor}
\newtheorem{thm}{Theorem}
\newtheorem{prop}{Proposition}
\newtheorem{claim}{Claim}
\newcommand{\NIP}{\mathrm{NIP}}
\newcommand{\U}{\mathcal{U}}
\newcommand{\Ualg}{\mathcal{U}_0}
\title{Residually dominated groups in henselian valued fields of equicharacteristic zero}
\author{Dicle Mutlu, Paul Z. Wang}
\date{Source: arXiv:2502.08530v2 (CC BY 4.0)}
\begin{document}
\maketitle

\noindent\emph{Source and licence.} Residually dominated groups in henselian valued fields of equicharacteristic zero. Version arXiv:2502.08530v2 (CC BY 4.0), distributed by arXiv under the Creative Commons Attribution 4.0 International licence, http://creativecommons.org/licenses/by/4.0/, as recorded in the arXiv licence metadata for this exact version and shown on the abstract page https://arxiv.org/abs/2502.08530v2. Full licence notice: \texttt{licenses/arxiv-2502.08530-CC-BY-4.0.txt}.\par\smallskip

\noindent\textit{Editorial scope.} Counted unit: the article's prop prop\_functoriality (source0.tex lines 1276--1286). The article's complete original proof of it is delivered with it (source0.tex lines 1288--1324, one \texttt{proof} environment of 37 lines). No mathematics is written, repaired or re-derived: the statement and the proof are the article's own lines, in the article's own notation, and no retained line is substituted or re-flowed.  Retained uncounted as the source-local context the proof needs: lines 1186--1201.  Nothing was omitted from any retained span, so the package carries no omission record.  No retained line contains a citation, so the wrapper carries no bibliography and no bibliography entry.  The retained lines contain the article's own diagram code, which the wrapper typesets with the standard tikz packages; no image file is delivered and the package declares no asset dependency.  Editorial formatting only, in addition: an AMS \texttt{amsart} preamble that restates the article's own declarations and macros used by the retained lines, namely the environments \texttt{thm}, \texttt{prop}, \texttt{claim} on separate counters, together with the standard packages those lines need.  The retained environments print in this document as Theorem 1, Proposition 1, Claim 1 in order of appearance, whereas the article prints the same items with the numbering of its own full document.  The complete native source of the article as distributed in the arXiv e-print for this exact version is bundled unchanged as \texttt{sources/173177/source0.tex}.
\begin{thm}\label{thm_acvf_enveloppe} Assume $T$ is $\NIP$.
    Let $(G,\cdot)$ be a residually dominated group lying in the valued field sort of $\U$. Then, there exists an $\mathcal{L}_0$-type-definable group $G_1$, a definable group homomorphism $\iota: G^{00} \rightarrow G_1$ with finite kernel, with the following properties:

    \begin{enumerate}
        \item The group $G_1$ is, in the structure $\Ualg$, connected stably dominated, and the image of any generic of $G^{00}$ is generic, in the sense of $ACVF$, in $G_1$.
        
        \item For all $\mathcal{L}_0$-type-definable groups $H_1$, for all definable morphisms $f: G^{00} \rightarrow H_1$, there exist $\mathcal{L}_0$-type-definable $H'_1 \leq H_1$, $H_0 \lhd H'_1$ finite, and $\psi: G_1 \rightarrow H'_1/H_0$ such that $f$ factors through $H'_1$ and $\pi \circ f = \psi \circ \iota$, where $\pi : H'_1 \rightarrow H'_1 / H_0$ is the quotient map. In other words, the following diagram commutes:
    \end{enumerate}

\begin{center}
\begin{tikzcd} G^{00}\arrow[d, "f"] \arrow[r, "\iota"]& G_1 \arrow[d, "\psi"] \\
 H'_1\arrow[r, "\pi"]& H'_1 / H_0  \end{tikzcd}
\end{center}

    Moreover, such a group $G_1$ is unique up to $\mathcal{L}_0$-definable isogeny. 
\end{thm}

\begin{prop}\label{prop_functoriality} Assume $T$ is $\NIP$. Let $(G, \cdot)$, $(H, \cdot)$ be residually dominated definable groups, in the valued field sort. Let $f: G \twoheadrightarrow H$ be a surjective definable group homomorphism. Let $G_1$, $H_1$ be $\mathcal{L}_0$-definable groups, and $\iota_G: G^{00} \rightarrow G_1$, $\iota_H: H^{00} \rightarrow H_1$ be definable group homomorphisms satisfying the conclusions of Theorem \ref{thm_acvf_enveloppe}.


Then, there exists a finite normal subgroup $H_0 \lhd H_1$ and an $\mathcal{L}_0$-definable surjective (in ACVF) group homomorphism $\psi: G_1 \twoheadrightarrow H_1 / H_0$ such that the following diagram commutes:

\begin{center}
\begin{tikzcd} G^{00}\arrow[d, "f"] \arrow[rr, "\iota_G"]&& G_1 \arrow[d, "\psi"] \\
 H^{00} \arrow[r, "\iota_H"]& H_1 \arrow[r] & H_1 / H_0  \end{tikzcd}
\end{center}

\end{prop}

\begin{proof}
\setcounter{claim}{0}
    Let $M$ be a sufficiently saturated model over which everything is defined. Let $p \in S_G(M)$ be a residually dominated generic concentrating on $G^{00}$.
    \begin{claim}
        The type $f_*(p) \in S_H(M)$ is a residually dominated generic of $H$ concentrating on $H^{00}$.
    \end{claim}
    \begin{proof}
        By $\NIP$, we know that $Stab_G(p)=G^{00}$, so $p$ concentrates on $Stab_G(p)$. Then, since $f$ is a group homomorphism, this implies that $f_*(p)$ concentrates on $Stab_H(f_*(p))$. Now, by surjectivity of $f$, and genericity of $p$, the type $f_*(p)$ is generic in $H$, so $Stab_H(f_*(p))=H^{00}$. Thus, the type $f_*(p)$ concentrates on $H^{00}$.

        To conclude, recall that residual domination is preserved under pushforward by definable maps. So $f_*(p)$ is residually dominated.
    \end{proof}

    Then, we apply the second point of Theorem \ref{thm_acvf_enveloppe}, to the composite $\iota_H \circ f: G^{00} \rightarrow H^{00} \rightarrow H_1$. So, there exists $\mathcal{L}_0$-definable subgroups $H'_1 \leq H_1$, $H_0 \lhd H'_1$ finite, and an $\mathcal{L}_0$-definable morphism $\psi: G_1 \rightarrow H'_1/H_0$ such that $\iota_H \circ f$ factors through $H'_1$ and $\pi \circ \iota_H \circ f = \psi \circ \iota_G$. In other words, the following diagram commutes:

    
\begin{center}
\begin{tikzcd} G^{00} \arrow[rd, "\iota_H \circ f"] \arrow[rr, "\iota_G"]&& G_1 \arrow[d, "\psi"] \\
 & H'_1 \arrow[r] & H'_1 / H_0  \end{tikzcd}
\end{center}

\begin{claim}
    The subgroup $H'_1 \leq H_1$ is equal to $H_1$.
\end{claim}
\begin{proof}
    Recall that the type $f_*(p)$ is residually dominated generic of $H^{00}$. Thus, by the first point of Theorem \ref{thm_acvf_enveloppe} applied to $H$, the $\mathcal{L}_0$-restriction of the type ${(\iota_H \circ f)}_*(p)$ is the unique generic of $H_1$ in $\Ualg$. Since $\iota_H \circ f$ factors through $H'_1$, this implies that ${(\iota_H \circ f)}_*(p)$ concentrates on $H'_1$. Since this type is generic in $H_1$, it generates $H_1$ in two steps, thus $H_1 \leq H'_1$, as required.
\end{proof}

Thus, we just proved that the following diagram makes sense, and commutes:

\begin{center}
\begin{tikzcd} G^{00} \arrow[d, "f"] \arrow[rd, "\iota_H \circ f"] \arrow[rr, "\iota_G"]&& G_1 \arrow[d, "\psi"] \\ H^{00} \arrow[r, "\iota_H"]
 & H_1 \arrow[r] & H_1 / H_0  \end{tikzcd}
\end{center}


It remains to prove that, in $\Ualg$, the morphism $\psi$ is surjective. Using commutativity of the outer square, and recalling that the $\mathcal{L}_0$-restriction of ${(\iota_H \circ f)}_*(p)$ is generic in $H_1$, we deduce that the image of $\psi$ contains an generic of $H_1 / H_0$ in $\Ualg$. By connectedness of the latter group in $\Ualg$, this implies that $\psi$ is surjective, as required.
\end{proof}

\end{document}
